% Generated draft: requires mathematical and source review.
\documentclass[11pt]{article}
\usepackage[margin=25mm]{geometry}
\usepackage{fontspec}
\setmainfont{Latin Modern Roman}
\usepackage{amsmath,amssymb,mathtools,unicode-math}
\setmathfont{Latin Modern Math}
\usepackage{enumitem,needspace,xurl,hyperref}
\setlength{\parindent}{0pt}
\setlength{\parskip}{5pt}
\setlength{\emergencystretch}{3em}
\newcommand{\N}{\mathbb{N}}
\newcommand{\Z}{\mathbb{Z}}
\newcommand{\Q}{\mathbb{Q}}
\newcommand{\R}{\mathbb{R}}
\newcommand{\C}{\mathbb{C}}
\newcommand{\F}{\mathbb{F}}
\newcommand{\A}{\mathbb{A}}
\newcommand{\PP}{\mathbb{P}}
\newcommand{\E}{\mathbb{E}}
\newcommand{\eps}{\varepsilon}
\newcommand{\dd}{\,\mathrm{d}}
\begin{document}
\section*{UnsolvedMath ID: 800263}
\textbf{Source problem code:} catalogue entry
\subsection{Definitions and mathematical statement}
Let \(n\geq 1\) be an integer and let \(F\in\C[w_0,\ldots,w_n]\). Set
\[
X_F=\{w\in\C^{n+1}:F(w)=1\}.
\]
Assume that this equation defines a smooth affine algebraic hypersurface of complex dimension \(n\); explicitly,
\[
\left(\frac{\partial F}{\partial w_0}(p),\ldots,\frac{\partial F}{\partial w_n}(p)\right)\neq 0
\qquad\text{for every }p\in X_F.
\]
Its complex tangent space at \(p\) is
\[
T_pX_F=\left\{v\in\C^{n+1}:\sum_{j=0}^{n}\frac{\partial F}{\partial w_j}(p)v_j=0\right\}.
\]
\textbf{Question (R. Narasimhan).} Does every such \(X_F\) admit a holomorphic immersion into \(\C^n\)? That is, does there exist a holomorphic map \(f:X_F\to\C^n\) such that
\[
\forall p\in X_F,\qquad \dd f_p:T_pX_F\longrightarrow\C^n
\quad\text{is injective}?
\]
Since both tangent spaces have complex dimension \(n\), this asks for an everywhere invertible differential; global injectivity is not required.

The source states parenthetically that these hypersurfaces are always parallelizable, attributing the case \(n=2\) to Murthy and Swan and general \(n\) to Suslin. Parallelizability means triviality of the tangent bundle in a specified category, but the source does not specify that category. This assertion is retained as supplied background, not as an additional hypothesis or an independently verified theorem. The complex coefficient field and positive-integer range of \(n\) make explicit the natural complex-geometric reading; the source does not state them separately.
\subsection{Short English statement}
Must every smooth affine hypersurface \(F=1\) in \(\C^{n+1}\) admit a holomorphic map to \(\C^n\) whose differential is injective everywhere? The source attributes this question to R. Narasimhan and supplies a parallelizability assertion with an unspecified category.
\subsection{Sources}
\begin{itemize}
\item Local UnsolvedMath registry, permanent ID \(800263\), exact source problem code \texttt{catalogue entry}; title: \emph{Several Complex Variables Problem List --- Problem 62}. The supplied record labels the problem open; that label is not an independently verified assessment of its present status.
\item Buzzard, Fornaess, Krantz, \emph{The Several Complex Variables Problem List} (1995), Problem 62. Supplied source URL: \url{https://arxiv.org/abs/math/9509202}. The record reports access on 2026-10-16 and extraction by \texttt{arxiv\_latex}; these are supplied provenance details, not a claim of source access for this draft.
\item The problem text attributes the question to R. Narasimhan, Department of Mathematics, University of Chicago, Chicago, IL 60637. The registry separately lists Buzzard, Fornaess, Krantz as proposers, with proposed year 1995.
\item The source attributes its parallelizability assertion to Murthy and Swan for \(n=2\), and to Suslin for general \(n\), without supplying publication titles or bibliographic references for those results.
\end{itemize}
\end{document}
